\chapter{Complexity — Parse2}
%%% generated by the blueprint pipeline from TCSlib.Complexity.SpaceComplexity.Machines.Parse2
\section{Overview}

Declarations of the Lean module \texttt{TCSlib.Complexity.SpaceComplexity.Machines.Parse2}.

\section{Declarations}

\begin{definition}[Scanning Boolean pairs]
\lean{Complexity.LogProg.scanPairs}
\label{Complexity.LogProg.scanPairs}
\leanok
Define a function that takes an optional previous Boolean value and a list of Boolean
values, and returns either no result or a list. For a list beginning with false, true,
it returns no result if the previous value is false, and otherwise returns the remaining
list; for a list beginning with two equal values, it recursively scans the remaining
list with that value as the previous value. It returns no result in every other case.
\end{definition}

\begin{lemma}[Last element after prefixing a list]
\lean{Complexity.LogProg.getLast?_cons_of}
\leanok
\difficulty{4}
Let $u$ be a Boolean list whose last-element operation returns $v$. Then prefixing $u$
by any Boolean $b$ leaves its last element equal to $v$.
\end{lemma}

\begin{lemma}[Characterization of successful pair scans]
\lean{Complexity.LogProg.scanPairs_iff}
\label{Complexity.LogProg.scanPairs_iff}
\leanok
\uses{Complexity.LogProg.getLast?_cons_of, Complexity.LogProg.scanPairs, Turing.dbl, Turing.dbl_cons, Turing.dbl_nil}
\difficulty{5}
For an initial remembered bit $\ell\in\{\mathrm{none},0,1\}$ and Boolean words $z,w$,
the pair scan started with $\ell$ returns $w$ exactly when there is a Boolean word $u$
such that $z=u_0u_0u_1u_1\cdots u_{n-1}u_{n-1}\,01w$ and the last bit of $u$, if
present, or otherwise $\ell$, is not $0$.
\end{lemma}

\begin{definition}[Rejecting computation with register-position constraints]
\lean{Complexity.LogProg.Rejects}
\label{Complexity.LogProg.Rejects}
\leanok
\uses{Complexity.LogProg.RProg, Complexity.LogProg.rrun, Complexity.LogProg.xCfg, Turing.Cfg}
Let $P$ be a register program, let an oracle assign Boolean answers to each decider and
Boolean-list input, and let $c_0,c$ be configurations on input $x$. For a register $r$
and integer position $p$, there must be a time $T$ such that the run from $c_0$ is
halted at time $T$, its output is the output of $c$ followed by $0$, and its work-tape
head positions agree with those of $c$ except that the head on register $r$ is at $p$.
At every earlier time $t<T$, the run must be in a configuration of the specified form
determined by $c$, some state and input-head position, $r$, and $p$, and that state must
not be a call node.
\end{definition}

\begin{definition}[Reachable computation with register-position constraints]
\lean{Complexity.LogProg.Reaches}
\label{Complexity.LogProg.Reaches}
\leanok
\uses{Complexity.LogProg.RProg, Complexity.LogProg.rrun, Complexity.LogProg.xCfg, Turing.Cfg}
Let $P$ be a register program, let an oracle assign Boolean answers to each decider and
Boolean-list input, and let $c_0,c_1,c$ be configurations on input $x$. For a register
$r$ and integer position $p$, there must be a time $T$ such that the run from $c_0$ at
time $T$ equals $c_1$. At every earlier time $t<T$, the run must be in a configuration
of the specified form determined by $c$, some state and input-head position, $r$, and
$p$, and that state must not be a call node.
\end{definition}

\begin{lemma}[One-step extension of reaching]
\lean{Complexity.LogProg.Reaches.step}
\label{Complexity.LogProg.Reaches.step}
\leanok
\uses{Complexity.LogProg.Reaches, Complexity.LogProg.rrun, Complexity.LogProg.rrun_add, Complexity.LogProg.xCfg, Turing.Cfg}
\difficulty{4}
Let $P$ be a register program with oracle, and let $c_0,c_1,c_2,c$ be configurations on
the same input, with register $r$, integer position $p$, state $s$, and input-head
position $ip$. Suppose $c_0$ is the configuration obtained from $c$ by setting its state
to $s$, its input-head position to $ip$, and register $r$'s head position to $p$;
suppose also that $P$ has no oracle call at $s$, that one step from $c_0$ yields $c_1$,
and that $c_1$ reaches $c_2$ relative to $c,r,p$. Then $c_0$ reaches $c_2$ relative to
$c,r,p$.
\end{lemma}

\begin{lemma}[One-step extension of rejection]
\lean{Complexity.LogProg.Rejects.step}
\label{Complexity.LogProg.Rejects.step}
\leanok
\uses{Complexity.LogProg.Rejects, Complexity.LogProg.rrun, Complexity.LogProg.rrun_add, Complexity.LogProg.xCfg, Turing.Cfg}
\difficulty{4}
Let $P$ be a register program with oracle, and let $c_0,c_1,c$ be configurations on the
same input, with register $r$, integer position $p$, state $s$, and input-head position
$ip$. Suppose $c_0$ is obtained from $c$ by setting its state to $s$, its input-head
position to $ip$, and register $r$'s head position to $p$; suppose that $P$ has no
oracle call at $s$, that one step from $c_0$ yields $c_1$, and that the computation from
$c_1$ rejects relative to $c,r,p$. Then the computation from $c_0$ rejects relative to
$c,r,p$.
\end{lemma}

\begin{lemma}[Transitivity of reaching]
\lean{Complexity.LogProg.Reaches.trans}
\label{Complexity.LogProg.Reaches.trans}
\leanok
\uses{Complexity.LogProg.Reaches, Complexity.LogProg.rrun_add, Turing.Cfg}
\difficulty{4}
Let $P$ be a register program with oracle, and let $c_0,c_1,c_2,c$ be configurations on
the same input, with register $r$ and integer position $p$. If $c_0$ reaches $c_1$
relative to $c,r,p$ and $c_1$ reaches $c_2$ relative to $c,r,p$, then $c_0$ reaches
$c_2$ relative to $c,r,p$.
\end{lemma}

\begin{lemma}[Composition of reaching and rejection]
\lean{Complexity.LogProg.Reaches.rejects}
\label{Complexity.LogProg.Reaches.rejects}
\leanok
\uses{Complexity.LogProg.Reaches, Complexity.LogProg.Rejects, Complexity.LogProg.rrun_add, Turing.Cfg}
\difficulty{4}
Let $P$ be a register program with oracle, and let $c_0,c_1,c$ be configurations on the
same input, with register $r$ and integer position $p$. If $c_0$ reaches $c_1$ relative
to $c,r,p$, and the computation from $c_1$ rejects relative to $c,r,p$, then the
computation from $c_0$ rejects relative to $c,r,p$.
\end{lemma}

\begin{lemma}[Immediate rejection at a non-call state]
\lean{Complexity.LogProg.rejects_one}
\label{Complexity.LogProg.rejects_one}
\leanok
\uses{Complexity.LogProg.Rejects, Complexity.LogProg.inSym, Complexity.LogProg.rejAct, Complexity.LogProg.rrun_one_rej, Complexity.LogProg.xCfg, Turing.Cfg}
\difficulty{4}
Let $P$ be a register program, let $c$ be a configuration, and let $r$ be a register
with head position $p$. Suppose state $s$ is not a call state, and for every possible
tuple of work-tape symbols the transition of the program's underlying machine at $s$, on
the input symbol read at position $ip$, is the action that emits false and halts without
moving any head or writing to any tape. Then, starting from the configuration obtained
from $c$ by setting the state to $s$, the input head to $ip$, and register $r$'s head to
$p$, the program rejects in one step in the sense that its halting configuration has
output $c$'s output followed by false, has the specified adjusted head-position map, and
all pre-final configurations are ordinary non-call program states.
\end{lemma}

\begin{lemma}[Reflexivity of reaching]
\lean{Complexity.LogProg.Reaches.refl}
\label{Complexity.LogProg.Reaches.refl}
\leanok
\uses{Complexity.LogProg.Reaches, Turing.Cfg}
\difficulty{2}
Let $P$ be a register program with oracle, let $c_0$ and $c$ be configurations on the
same input, let $r$ be a work-tape register, and let $p\in\mathbb{Z}$. Then $c_0$
reaches itself relative to the reference configuration $c$, register $r$, and position
$p$.
\end{lemma}

\begin{lemma}[One-step rightward head movement]
\lean{Complexity.LogProg.reaches_right}
\label{Complexity.LogProg.reaches_right}
\leanok
\uses{Complexity.LogProg.apply_xAct, Complexity.LogProg.inSym, Complexity.LogProg.moveInputPos_pos_val, Complexity.LogProg.rrun, Complexity.LogProg.rrun_one_x, Complexity.LogProg.xAct, Complexity.LogProg.xCfg, Turing.Cfg}
\difficulty{3}
Let $P$ be a register program with $m$ work tapes, input $x$, oracle $O$, and register
$r$. Suppose state $s$ is not a call state, and for every tuple of work-tape symbols its
transition in state $s$, when the input symbol at position $q$ is read, moves the input
head one position to the right, leaves all work-tape contents and head positions
unchanged, emits no output, and enters state $s'$. For $q\le |x|$, any initial
configuration with state $s$, input-head position $q$, and work-tape head $r$ at
position $p$ reaches in one program step the same configuration data with state $s'$ and
input-head position $q+1$.
\end{lemma}

\begin{lemma}[Characterization of canonical words]
\lean{Complexity.LogProg.canon_iff}
\label{Complexity.LogProg.canon_iff}
\leanok
\uses{Complexity.LogProg.Canon}
\difficulty{4}
A Boolean word $u$ is canonical if and only if its last-bit value, when defined, is not
false.
\end{lemma}

\begin{lemma}[Left identity for option choice]
\lean{Complexity.LogProg.none_orElse'}
\label{Complexity.LogProg.none_orElse'}
\leanok
\difficulty{1}
For every optional Boolean value $o$, choosing the first available value between no
value and $o$ yields $o$.
\end{lemma}

\begin{definition}[State selected by an optional Boolean]
\lean{Complexity.LogProg.wSt}
\label{Complexity.LogProg.wSt}
\leanok
Given three states $v_0,v_F,v_T$ and an optional Boolean, define a state-valued function
that returns $v_0$ when the Boolean is absent, $v_F$ when it is false, and $v_T$ when it
is true.
\end{definition}

\begin{lemma}[Simulation of a Boolean word scan]
\lean{Complexity.LogProg.valWord_run}
\label{Complexity.LogProg.valWord_run}
\leanok
\uses{Complexity.LogProg.Reaches, Complexity.LogProg.Reaches.step, Complexity.LogProg.Rejects, Complexity.LogProg.Rejects.step, Complexity.LogProg.apply_xAct, Complexity.LogProg.getLast?_cons_of, Complexity.LogProg.inSym, Complexity.LogProg.none_orElse', Complexity.LogProg.reaches_right, Complexity.LogProg.rejects_one, Complexity.LogProg.rewind_x, Complexity.LogProg.rrun, Complexity.LogProg.rrun_one_x, Complexity.LogProg.valWAct, Complexity.LogProg.wSt, Complexity.LogProg.xAct, Complexity.LogProg.xCfg, Turing.Cfg}
\difficulty{6}
Let $P$ be a register program with oracle, let $r$ be a register, and let
$v_0,v_F,v_T,r_1,r_2,n$ be states. Suppose the transitions from $v_0,v_F,v_T$ act
according to the value word action, with the remembered last bit respectively absent,
false, and true; suppose $r_1$ moves the input head left and enters $r_2$, while $r_2$
moves left and stays in $r_2$ when it reads a symbol, and moves right and enters $n$ at
the right boundary. Assume none of these five states is a call state. For a
configuration $c$, integer position $p$, Boolean word $x$, word $w$, and natural number
$q$ satisfying $q+|w|=|x|+1$, assume that at every offset $k$ the symbol read at input
position $q+k$ is the $k$-th bit of $w$. Starting at input position $q$ in state
$v_0,v_F$, or $v_T$ according as the remembered bit is absent, false, or true, with
register $r$'s head at $p$ and all other configuration data inherited from $c$, the run
reaches the analogous configuration in state $n$ at input position $1$ if the last bit
of $w$ (or the remembered bit when $w$ is empty) is not false; if that bit is false, it
instead halts after appending a rejecting output bit, with register-head positions equal
to those of $c$ except that register $r$ is at $p$. In either case, every earlier
configuration is in a state with no call and has the same register-head positions, with
register $r$ at $p$.
\end{lemma}

\begin{definition}[Action for the first verifier]
\lean{Complexity.LogProg.valP1Act}
\label{Complexity.LogProg.valP1Act}
\leanok
\uses{Complexity.LogProg.rejAct, Complexity.LogProg.xAct, Turing.Action}
Given a register index $r$, a function assigning a state to each Boolean, and an
optional Boolean answer, define a machine action as follows: for answer $b$, use the
action prescribed by $r$ and the state assigned to $b$; if the answer is absent, use the
rejection action.
\end{definition}

\begin{definition}[Action for the second verifier]
\lean{Complexity.LogProg.valP2Act}
\label{Complexity.LogProg.valP2Act}
\leanok
\uses{Complexity.LogProg.rejAct, Complexity.LogProg.xAct, Turing.Action}
Given a register index $r$, states $v_F,v_T,v_0$, a previous optional Boolean value
$\text{last}$, a Boolean $b_1$, and an optional Boolean answer, define an action as
follows. If the answer is absent, use the rejection action. If it is $b_2$, then when
$b_1$ is false and $b_2$ is true, use the rejection action if $\text{last}$ is false, and
otherwise use the action prescribed by $r$ and $v_0$; when $b_1=b_2$, use the action
prescribed by $r$ and $v_T$ if $b_1$ is true and by $r$ and $v_F$ if $b_1$ is false; in
the remaining case use the rejection action.
\end{definition}

\begin{lemma}[Scanning paired encodings]
\lean{Complexity.LogProg.valPairs_run}
\label{Complexity.LogProg.valPairs_run}
\leanok
\uses{Complexity.LogProg.Reaches, Complexity.LogProg.Reaches.refl, Complexity.LogProg.Reaches.step, Complexity.LogProg.Rejects, Complexity.LogProg.Rejects.step, Complexity.LogProg.inSym, Complexity.LogProg.reaches_right, Complexity.LogProg.rejects_one, Complexity.LogProg.rrun, Complexity.LogProg.scanPairs, Complexity.LogProg.valP2Act, Complexity.LogProg.xCfg, Turing.Cfg}
\difficulty{6}
Let $P$ be a register program with oracle $\mathcal O$, $r$ a register, $p\in\mathbb Z$,
and $c$ a configuration on input $x$. Suppose the transition identities $hP1,hP2$ hold
at states $vP1,vP2$ and neither state is an oracle-call state, as asserted by $cP1,cP2$.
For every Boolean word $z$, optional last bit $last$, and natural number $q$ satisfying
$q+|z|=|x|+1$, assume that for each $k$, the symbol of $x$ at one-based position $q+k$
equals the $k$th bit of $z$, when defined. If scanning the aligned pairs of $z$, with
initial previous bit $last$, returns a word $w$, then the run from state $vP1(last)$ at
input-head position $q$ reaches state $vW0$ at position $q+|z|-|w|$, which is a valid
input-head position. If the scan returns no word, that run halts with a rejecting
output. In both cases, the stated reachability or rejection runs avoid oracle-call
states before their endpoint.
\end{lemma}

\begin{lemma}[Validation of the unary prefix]
\lean{Complexity.LogProg.valPrefix_run}
\label{Complexity.LogProg.valPrefix_run}
\leanok
\uses{Complexity.LogProg.Reaches, Complexity.LogProg.Reaches.refl, Complexity.LogProg.Reaches.rejects, Complexity.LogProg.Reaches.step, Complexity.LogProg.Reaches.trans, Complexity.LogProg.Rejects, Complexity.LogProg.getElem?_lt_tRun, Complexity.LogProg.getElem?_tRun, Complexity.LogProg.inSym, Complexity.LogProg.inSym_succ, Complexity.LogProg.reaches_right, Complexity.LogProg.rejects_one, Complexity.LogProg.rrun, Complexity.LogProg.scanR, Complexity.LogProg.tRun, Complexity.LogProg.tRun_le, Complexity.LogProg.valUAct, Complexity.LogProg.xCfg, Turing.Cfg}
\difficulty{6}
Let $P$ be a register program with oracle $\mathcal O$, $r$ a register, $p\in\mathbb Z$,
and $c$ a configuration on input $x$. Suppose the transition identities $hU0,hU1,hS$
hold at states $vU0,vU1,vS$, respectively, and none of these states is an oracle-call
state, as asserted by $cU0,cU1,cS$. If the length of the initial run of ones in $x$ is
even and the next two bits are $01$, then the run from state $vU0$ at input-head
position $1$ and register-head position $p$ reaches state $vX$ at input-head position
$3+tRun(x)$, a valid position. If this bit-pattern condition fails, the run from the
same starting configuration halts with a rejecting output. The runs avoid oracle-call
states before their endpoint.
\end{lemma}

\begin{lemma}[Pairing decomposition at a prefix separator]
\lean{Complexity.LogProg.prefix_decomp}
\label{Complexity.LogProg.prefix_decomp}
\leanok
\uses{Complexity.LogProg.dbl_replicate, Complexity.LogProg.getElem?_lt_tRun, Complexity.LogProg.tRun, Complexity.LogProg.tRun_le, Turing.pairEncode, Turing.pairEncode_eq_dbl}
\difficulty{4}
Let $y$ be a bit string, and let $r$ be the length of its initial run of ones. If $r$ is
even, the bit of $y$ at index $r$ is $0$, and the bit at index $r+1$ is $1$, then $y$ is
the pairing of the unary string of $r/2$ ones with the suffix of $y$ beginning at index
$r+2$.
\end{lemma}

\begin{lemma}[Initial run in a paired encoding]
\lean{Complexity.LogProg.tRun_pairEncode}
\label{Complexity.LogProg.tRun_pairEncode}
\leanok
\uses{Complexity.LogProg.dbl_replicate, Complexity.LogProg.tRun, Turing.pairEncode, Turing.pairEncode_eq_dbl}
\difficulty{3}
For every natural number $n$ and Boolean list $z$, the initial run of true entries in
the pairing of the length-$n$ all-true list with $z$ has length $2n$.
\end{lemma}

\begin{lemma}[Characterization of valid paired words]
\lean{Complexity.LogProg.validPair_iff}
\label{Complexity.LogProg.validPair_iff}
\leanok
\uses{Complexity.LogProg.Canon, Complexity.LogProg.ValidPair, Complexity.LogProg.canon_iff, Complexity.LogProg.dbl_replicate, Complexity.LogProg.none_orElse', Complexity.LogProg.prefix_decomp, Complexity.LogProg.scanPairs, Complexity.LogProg.scanPairs_iff, Complexity.LogProg.tRun, Complexity.LogProg.tRun_pairEncode, Turing.dbl, Turing.pairEncode, Turing.pairEncode_eq_dbl}
\difficulty{4}
For every Boolean word $y$, validity is equivalent to both of the following: the length
of its initial run of ones is even, and the next two bits are $01$; moreover, scanning
aligned pairs in the suffix after those bits, starting with no previous pair, succeeds
with some remaining word $w$, and $w$ is canonical (empty or ending in $1$).
\end{lemma}

\begin{lemma}[Validation of a paired input]
\lean{Complexity.LogProg.valPair_run}
\label{Complexity.LogProg.valPair_run}
\leanok
\uses{Complexity.LogProg.Canon, Complexity.LogProg.Reaches, Complexity.LogProg.Reaches.rejects, Complexity.LogProg.Reaches.trans, Complexity.LogProg.Rejects, Complexity.LogProg.ValidPair, Complexity.LogProg.canon_iff, Complexity.LogProg.inSym, Complexity.LogProg.inSym_succ, Complexity.LogProg.scanPairs, Complexity.LogProg.scanPairs_iff, Complexity.LogProg.tRun, Complexity.LogProg.valPairs_run, Complexity.LogProg.valPrefix_run, Complexity.LogProg.valWord_run, Complexity.LogProg.validPair_iff, Complexity.LogProg.wSt, Complexity.LogProg.xCfg, Turing.Cfg, Turing.dbl}
\difficulty{6}
Let $P$ be a register program with oracle $\mathcal O$, $r$ a register, $p\in\mathbb Z$,
and $c$ a machine configuration on input $x$. Assume the program's transitions at the
specified states implement the prescribed validation procedure: the states
$vU0,vU1,vS,vP1,vP2,vW0,vWF,vWT,rw_1,rw_2$ have the transition identities
$hU0,hU1,hS,hP1,hP2,hW0,hWF,hWT,h_1,h_2$, respectively, and none of these states is an
oracle-call state, as expressed by $cU0,cU1,cS,cP1,cP2,cW0,cWF,cWT,c_1,c_2$. Then, if
$x$ is a valid paired encoding, the run from the configuration based on $c$ at state
$vU0$, input-head position $1$, and register-head position $p$ reaches the corresponding
configuration at state $next$, input-head position $1$, and register-head position $p$.
If $x$ is not a valid paired encoding, that run instead halts with a rejecting output,
preserving the indicated register-head positions and avoiding oracle-call states before
halting.
\end{lemma}
